% Part: lambda-calculus
% Chapter: introduction
% Section: church-rosser

\documentclass[../../../include/open-logic-section]{subfiles}

\begin{document}

\olfileid[hi]{lam}{int}{cr}
\olsection{चर्च--रॉसर गुण}

\begin{thm}
\ollabel{thm:church-rosser}
$M$, $N_1$ और $N_2$ ऐसे पद हों कि $M \red N_1$ और $M \red N_2$। तब कोई
ऐसा पद $P$ है कि $N_1 \red P$ और $N_2 \red P$।
\end{thm}

\begin{cor}
मान लें $M$ को प्रसामान्य रूप में अपचयित किया जा सकता है। तब यह प्रसामान्य
रूप अद्वितीय है।
\end{cor}

\begin{proof}
यदि $M \red N_1$ और $M \red N_2$, तो पिछले प्रमेय से कोई ऐसा पद~$P$ है
कि $N_1$ और $N_2$ दोनों अपचयित होकर~$P$ बनते हैं। यदि $N_1$ और~$N_2$
दोनों प्रसामान्य रूप में हों, तो ऐसा केवल तभी हो सकता है जब $N_1 \ident P
\ident N_2$।
\end{proof}

अंततः, हम दो पदों $M$ और~$N$ को \emph{$\beta$-तुल्य}, अथवा केवल
\emph{तुल्य}, तब कहेंगे जब वे अपचयित होकर किसी समान पद में पहुँचते हों;
दूसरे शब्दों में, जब कोई ऐसा $P$ हो कि $M \red P$ और $N \red P$। इसे $M
\equal[\beta] N$ लिखते हैं। \olref{thm:church-rosser} का उपयोग करके आप जाँच
सकते हैं कि $\equal[\beta]$ एक तुल्यता-संबंध है, और उसमें यह अतिरिक्त गुण
है कि प्रत्येक $M$ और~$N$ के लिए, यदि $M \red N$ अथवा $N \red M$, तो $M
\equal[\beta] N$। (वास्तव में, यह दिखाया जा सकता है कि $\equal[\beta]$ इस
गुण वाला \emph{सबसे छोटा} तुल्यता-संबंध है।)

\end{document}

